\section{结果：接口为何需要证据来源约束}

结果分为两部分。首先，代表性记录说明测量侧能够把长时、多通道原始数据压缩为可回指的观测包，并以紧凑告警字段形成明确的检索意图。其次，受控扩库结果表明，在相关补充资料加入后，语义近邻得分与核心权威来源的可见性可以发生背离；查询表述、语言、编码器与分块方式会改变这种背离的幅度。本文据此讨论检索证据的可见性，而不将任何信号事件或候选解释提升为工程确认。

\subsection{代表性 trace：从多通道记录到受控检索意图}

以 A2 组的一条记录为例，初始化步骤生成了包含 597,607 个时间样本和 7 个中子电流通道的活跃记录；时间顺序校正后，记录覆盖约 7 d、名义采样间隔为 1 s，且未发现缺失值。确定性工具随后将事件、零值、通道关系和极值等结果写入记事板。该实例中，7 个通道各记录到 12 个孤立零值观测，短时偏离事件数为 30--34 个，通道平均值为 32.4 个。多个通道的零值时间重复出现，因此“同步零值—信号完整性/共享测量上下文”被保留为紧凑告警字段中的检索线索，而非可直接命名的故障类型。

这一追溯链的关键不在于事件数本身，而在于每个摘要字段均可回指：零值时间、事件计数和持续时间来自事件扫描输出；通道间共变、极值位置和快照来自相应的确定性工具；状况报告只组织这些已保存字段。资料检索后的候选说明不被提升为已确认根因。

据此，本案例的检索意图是查找与测量完整性、同步事件、冗余或共享仪表上下文有关的可定位技术段落。后续扩库实验表明，仅以语义近邻选择这些段落，不能保证预设核心资料仍在结果中可见。


\subsection{字符分块扩库协议：语义相似度与核心资料在前 10 个结果中的占比可以背离}

\texttt{nomic-embed-text} 字符分块扩库协议显示，补充资料的加入不必降低最高余弦相似度，却可能迅速降低核心来源在前 10 个结果中的占比。每个非零扩库规模在每个分块配置下均基于 100 次补充资料随机抽样取平均，随后再聚合 17 个已保存配置；下列 $\pm$ 值表示配置间的标准差。中文测量短语查询的核心资料前十名占比从仅含核心资料时的 1.00，在加入 5 份补充资料后降至 $0.10\pm0.09$，加入 10 份后降至 0；与此同时，平均最高余弦相似度由 0.55 上升至 0.63。英文查询的下降较缓：加入 5 份补充资料后核心资料前十名占比为 $0.86\pm0.05$，在加入 45 份后为 $0.38\pm0.05$，最高余弦相似度则维持在约 0.69--0.70。

\begin{figure}[H]
  \centering
  \begin{minipage}[t]{0.48\columnwidth}
    \centering\textbf{(a)}\par\vspace{1mm}
    \includegraphics[width=\linewidth]{fig04a_scope_priority}
  \end{minipage}\hfill
  \begin{minipage}[t]{0.48\columnwidth}
    \centering\textbf{(b)}\par\vspace{1mm}
    \includegraphics[width=\linewidth]{fig04b_top1_cosine}
  \end{minipage}
  \caption{字符分块扩库下的语义相似度与核心资料在前 10 个结果中的占比（权威资料可见性）。结果聚合 17 个已保存分块配置；阴影表示配置间标准差。(a) 核心资料前十名占比可随相关补充资料加入而下降；(b) 在相同协议中，平均最高余弦相似度提高或保持稳定。两者可同时发生。}
  \label{fig:scope-priority}
\end{figure}

\begin{figure}[!b]
  \centering
  \includegraphics[width=0.90\columnwidth]{fig05_nomic_loss_landscapes}
  \caption{全数据库条件下 \texttt{nomic-embed-text} 的中英文平均语义匹配程度。每个曲面为同一语言两条固定查询的平均值；横轴为分块长度，纵轴为重叠长度，高度为前 10 个结果的平均余弦相似度。箭头标出各语言的最高点。}
  \label{fig:nomic-loss-landscape}
\end{figure}

图~\ref{fig:scope-priority} 表明，相关补充资料可使语义相似度提高或保持稳定，同时挤出前 10 个结果中的核心资料；相似度因此不能单独作为证据来源约束的依据。

图~\ref{fig:nomic-loss-landscape} 固定全数据库条件，考察语言与字符分块对平均匹配度的影响。英文查询的平均余弦相似度为 0.665--0.674，高于中文的 0.635--0.647；但两张曲面均较平缓，分块--重叠组合带来的变化分别不超过 0.009 和 0.012，峰值分别位于英文 800/120 和中文 800/150。语言和分块会影响绝对分数，却不足以给出跨语言的通用最优配置；实际选择仍须结合图~\ref{fig:scope-priority} 的权威资料可见性及返回段落的相关性核查。

\subsection{MiniLM encoder 对照：query、语言与分块的交互}

% Keep this in-text evidence figure attached to the opening 5.3 result,
% with the same compact vertical rhythm as Fig.~\ref{fig:scope-priority}.
{\setlength{\intextsep}{5pt plus 1pt minus 1pt}
\setlength{\abovecaptionskip}{4pt}
\begin{figure}[H]
  \centering
  \begin{minipage}[t]{0.48\columnwidth}
    \centering\textbf{(a) English queries}\par\vspace{1mm}
    \includegraphics[width=\linewidth]{fig05a_query_effect_en}
  \end{minipage}\hfill
  \begin{minipage}[t]{0.48\columnwidth}
    \centering\textbf{(b) Chinese queries}\par\vspace{1mm}
    \includegraphics[width=\linewidth]{fig05b_query_effect_zh}
  \end{minipage}
  \caption{多语 MiniLM 编码器在加入 10 份补充资料时，技术细节查询相对基线查询的平均核心资料前十名占比变化。每点对应一个字符分块设置，左图纵轴标出两图共享的分块长度/重叠；正值表示技术细节查询家族的权威资料可见性较高。}
  \label{fig:query-effect}
\end{figure}
}

为避免将中文查询与英文向编码器的语言失配混入主比较，图~\ref{fig:query-effect} 仅使用多语 MiniLM 编码器，在三种字符分块设置下比较简短基线查询与由观测包线索形成的技术细节查询。每个点表示固定分块设置下，技术细节查询家族相对基线查询家族的平均权威资料可见性变化：英文为 $+0.05$ 至 $+0.14$，中文为 $-0.23$ 至 $+0.02$。在两种 MiniLM 编码器与三种字符分块构成的六种配置上，这一不对称性仍然存在：英文技术细节查询的平均核心资料前十名占比为 0.65，高于英文基线查询的 0.60；中文技术细节查询则为 0.52，低于中文基线查询的 0.62。完整扩库轨迹均从仅核心资料条件下的 1.00 出发，但下降速度、交叉位置及其对查询细节的响应随语言和配置而变。因此，观测包线索形成的技术细节 query 并不必然提高权威资料可见性；图~\ref{fig:configuration-interaction} 进一步以热图展开这一条件性交互。

\begin{figure}[H]
  \centering
  \includegraphics[width=0.92\columnwidth]{fig06_configuration_heatmap}
  \caption{$m=10$ 时查询、语言、编码器与字符分块方式的交互。每个单元格为核心资料前十名占比的平均值；横向分隔线区分英文与中文查询，纵向分隔线区分两种编码器。由观测包线索形成的技术细节查询并不在所有语言和配置下提高权威资料可见性。}
  \label{fig:configuration-interaction}
\end{figure}

图~\ref{fig:configuration-interaction} 揭示了图~\ref{fig:query-effect} 中家族平均值所掩盖的条件性交互。图中数值均为预设核心资料在前 10 个结果中的占比，即本文定义的权威资料可见性。以“高同步”技术细节查询为例，MiniLM-L12 在英文三个分块设置下为 0.83--0.93，对应中文仅为 0.49--0.56；同一观测包线索并未在两种语言中对称提高权威资料可见性。图中也可见，更短、更宽泛的核心术语 query 有时取得较高占比，峰值位置则随语言、编码器和分块改变。因此，图7不支持由 encoder 名称、层数或跨配置平均值直接推出统一最优分块；部署时应先固定目标语言与 query 表述，再在相应 encoder--分块组合内按权威资料可见性筛选，并复核返回段落的相关性。

MiniLM encoder 对照与 \texttt{nomic-embed-text} 字符分块扩库协议共同表明：在本研究固定的资料库、检索规则和已测试 encoder 范围内，相关补充资料扩展后权威资料可见性下降的风险并非单一表示模型所致。由于两类协议的 query 集和部分分块设置不同，本文只比较各自协议内 CorePriority@10 随扩库规模的变化，不据此横向排列 encoder 或绝对余弦分数。技术细节 query 的作用随语言和配置而变，因此紧凑告警字段应被理解为受控的检索表述，而非必然提高权威资料可见性的查询扩展。

\subsection{分块边界与检索单元核查}

固定词数分块的 MiniLM 对照表明，上述风险并非只出现在 \texttt{nomic-embed-text} 的字符分块中。多语种 MiniLM 对中文短语产生的最高余弦相似度约为 0.62--0.64，高于另一 MiniLM 的约 0.27；但随着补充资料加入，两种 encoder 的 CorePriority@10 均出现下降。语义表示对不同语言的近邻得分具有实质影响，却不能替代来源要求。

\texttt{nomic-embed-text} 索引上的自检索则说明，检索单元的定义也会影响“找回”的含义。图~\ref{fig:self-retrieval} 选取 800/50、1000/80 与 1200/150 三个配置：800/50 是此前使用的基准；在完整的 $4\times4$ 网格中，1000/80 在加入 45 份补充资料时两种匹配标准的平均倒数排名差距最小（0.026），1200/150 的差距最大（0.157）。因此，后二者用于展示边界效应由弱到强的两个端点，而非宣称配置优劣。

图中的空心点是仅核心资料时的参考基线；非零扩库点则是补充资料随机抽样后的聚合结果。两者的抽样结构不同，故基线与第一个扩库点之间以虚线连接，不应将图首端的张开或收拢解释为“加入少量资料”带来的确定性性能跃迁。对每个非零扩库规模，两条实线之间的距离才表示严格同分块与同文档匹配在相同资料范围下的差异：原始段落未必以同一节点返回，但同一文档中的相邻分块可能已包含对应信息。因此，自检索结果用于解释分块边界与段落定位，而不应与权威资料可见性混为同一种证据。

\begin{figure}[!t]
  \centering
  \includegraphics[width=0.84\columnwidth]{fig08_self_retrieval_selected}
  \caption{三个代表性 \texttt{nomic-embed-text} 分块配置下的自检索结果。蓝线为严格同分块匹配，橙线为同文档匹配，阴影表示两种匹配标准的差距；纵轴为平均倒数排名，越高表示目标段落首次出现的位置越靠前。空心点为仅核心资料基线，虚线只连接基线与首个扩库点，实线比较相同扩库规模下的两种匹配标准。}
  \label{fig:self-retrieval}
\end{figure}

\FloatBarrier
\subsection{小结}

代表性 trace 说明，测量侧能够将无标签记录转换为可回指的检索线索；\texttt{nomic-embed-text} 字符分块扩库协议说明，证据侧若只依赖语义近邻，核心权威来源的可见性可能随相关资料的加入而丢失。MiniLM encoder 对照进一步表明：在本研究已测试的 encoder 范围内，这一风险并非单一向量表示所致；其幅度仍受 query、语言和分块方式共同影响，不能化约为某个组件的单一最优选择。与此同时，扩库后的权威资料可见性为检索协议选择提供实证依据。由此，来源 metadata 与显式来源策略不是检索实现的附属信息，而是 measurement-to-evidence interface 中需要被设计和评价的部分。
